\section{Données de l'IHM et du pupitre opérateur (contrôleur M221)}
\label{secDonneesIhm}

Le contrôleur M221 est scruté, comme la baie CS9 et comme le PFC200, par
un service d'IO-Scanning Ethernet/IP -- ici sur le réseau
\texttt{NOC\_ETHIP\_ROBOT}. Il assure à la fois l'îlot d'entrées/sorties du
pupitre opérateur (boutons, sélecteurs, voyants) et l'interface de
supervision Modbus-TCP vers l'écran HMI. Il produit 4 octets et consomme
40 octets. Les échanges sont structurés autour d'un objet de type
\texttt{THmi}.

\UPSTIremarque[Le M221 est lui aussi réservé en mémoire]{
Comme pour le PFC200 et la baie CS9, les données produites et consommées
par le M221 transitent par le service d'IO-Scanning Ethernet/IP : elles
occupent donc, elles aussi, une zone réservée de l'espace mémoire
\texttt{\%MW} de l'automate M340, à une adresse de base qui dépend de la
cartographie retenue pour le coupleur \texttt{NOC\_ETHIP\_ROBOT}. Cette
adresse n'est pas donnée ici : sa détermination (cf. document
\texttt{02c\_configuration\_noc}) fait partie du travail demandé aux
étudiants.
}

\subsection{Structure \texttt{THmi}}

\begin{center}
\begin{tabular}{|l|l|}
\hline
\textbf{Champ} & \textbf{Type} \\
\hline
\texttt{m\_Inputs} & \texttt{THmiInputs} (données produites -- éléments de
commande du pupitre) \\
\hline
\texttt{m\_Outputs} & \texttt{THmiOutputs} (données consommées -- voyants
du pupitre et écran de supervision) \\
\hline
\end{tabular}
\end{center}

\subsection{Données produites -- \texttt{THmiInputs} (éléments de commande)}

\begin{center}
\begin{tabular}{|l|p{8.5cm}|}
\hline
\textbf{Champ} & \textbf{Élément de commande} \\
\hline
\texttt{m\_nixJ1Plus} / \texttt{m\_nixJ1Minus} & Sélecteur de commande de
l'axe J1 (sens $+$ / sens $-$) \\
\hline
\texttt{m\_nixJ2Plus} / \texttt{m\_nixJ2Minus} & Sélecteur de commande de
l'axe J2 (sens $+$ / sens $-$) \\
\hline
\texttt{m\_nixJ3Up} / \texttt{m\_nixJ3Down} & Bouton-poussoir de commande
de l'axe J3 (montée / descente) \\
\hline
\texttt{m\_nixJ4Plus} / \texttt{m\_nixJ4Minus} & Sélecteur de commande de
l'axe J4 (sens $+$ / sens $-$) \\
\hline
\texttt{m\_nixAuto} / \texttt{m\_nixManu} & Sélecteur du mode de marche
souhaité (automatique / manuel) \\
\hline
\texttt{m\_nixEmptyTanks} & Bouton-poussoir « cuves vides » (confirmation
opérateur) \\
\hline
\texttt{m\_nixLoadTank0} / \texttt{m\_nixUnloadTank4} & Boutons-poussoirs
de chargement/déchargement manuel des cuves \\
\hline
\texttt{m\_nixEmptyTank0} / \texttt{m\_nixEmptyTank4} & Capteurs (contact
à ouverture) de cuve vide \\
\hline
\texttt{m\_nixStopRequest} & Bouton-poussoir « arrêt demandé » (contact à
ouverture) \\
\hline
\end{tabular}
\end{center}

\subsection{Données consommées -- \texttt{THmiOutputs}}

\texttt{THmiOutputs} se compose de deux sous-champs : \texttt{m\_qToR}
(voyants et affichage du pupitre) et \texttt{m\_Screen} (informations de
supervision transmises à l'écran HMI par Modbus-TCP).

\subsubsection{Sous-champ \texttt{m\_qToR} (voyants du pupitre)}

\begin{center}
\begin{tabular}{|l|p{8.5cm}|}
\hline
\textbf{Champ} & \textbf{Signification} \\
\hline
\texttt{m\_nqxRedColumn} / \texttt{m\_nqxGreenColumn} / \texttt{m\_nqxBlueColumn} &
Composantes rouge, verte, bleue de la colonne lumineuse du pupitre \\
\hline
\texttt{m\_nqxArmUnderPower} & Voyant « bras sous puissance » \\
\hline
\texttt{m\_na4qxHex7Seg} & Affichage 7 segments (ARRAY[0..3] OF BOOL,
codage binaire, [0]:LSB, [3]:MSB) \\
\hline
\texttt{m\_nqxLightLoadTank0..3} & Voyants de chargement associés à chacune
des quatre cuves \\
\hline
\texttt{m\_nqxLightUnLoadTank4} & Voyant de déchargement \\
\hline
\texttt{m\_nqxLightStopRequest} & Voyant « arrêt demandé » \\
\hline
\end{tabular}
\end{center}

\subsubsection{Sous-champ \texttt{m\_Screen} (supervision écran HMI)}

\begin{center}
\begin{tabular}{|l|l|p{6cm}|}
\hline
\textbf{Champ} & \textbf{Type} & \textbf{Signification} \\
\hline
\texttt{m\_iStateGmma} & INT & Numéro du mode GEMMA courant (affiché à
l'opérateur) \\
\hline
\texttt{m\_iCurvSpeed\_mm\_per\_s} & INT & Vitesse courante du robot,
exprimée en mm/s \\
\hline
\texttt{m\_a5tankLine} & ARRAY[0..4] OF \texttt{TTank} & État de chacune
des cinq lignes de traitement affichées à l'écran \\
\hline
\end{tabular}
\end{center}

Chaque élément de \texttt{m\_a5tankLine} est de type \texttt{TTank} :

\begin{center}
\begin{tabular}{|l|l|p{6.5cm}|}
\hline
\textbf{Champ} & \textbf{Type} & \textbf{Signification} \\
\hline
\texttt{m\_intPartNumber} & INT & Numéro identifiant la pièce en cours de
traitement \\
\hline
\texttt{m\_xBusy} & BOOL & Indique si la cuve est occupée \\
\hline
\texttt{m\_tElapsedTime} & TIME & Durée écoulée depuis le début du
traitement \\
\hline
\end{tabular}
\end{center}

\UPSTIremarque[Rôle de supervision]{
Ces informations ne participent à aucune commande directe de la partie
opérative : elles renseignent l'opérateur, via l'écran HMI, sur l'état
d'avancement du traitement dans chacune des cuves ainsi que sur le mode
GEMMA actif.
}
